\chapter{宽格子、集合散布与时域延长}
\label{ch:wide}

S1 每个格子 30 个种子，用来估计中位。分位是否被这 30 个种子钉住，要靠一张更宽的种子表来看。同一核的成员会不会自己散开，要靠一次没有分析步的集合积分。50 mm 在 55 min 内到不了设计线，究竟是时域不够还是通量不够，要把时域延长到 70 min 才分得清。厨房里对应的是：同一块厚肉糜，多煮一刻钟会不会跨过设计线，还是四个旋钮本身就送不进足够的热。下面三节按这个顺序写。

家用面板上的“智能菜单”常把同一厚度档写成一个确定的分钟数。宽格子把每个因子组合放到 30000 个种子上，只保留分位。15 mm 的带贴着上沿 \(+0.150\)，30 mm 跨过零点，50 mm 全部在零下、上沿仍低于 \(-0.3\)。带的宽度来自种子，厚度把整条带平移。把 30 个种子换成更宽的种子集，厚度的符号关系保持。货架彩页常配一张切开图：表面已经上色，中心不再是粉色。那张照片对应的是 30 mm、初温偏高、对偶组合这一条带的中段。同一菜单切到 50 mm，上沿仍低于 \(-0.3\)。分钟数写在厚度档上，内部热阻从未进入后验时，分位带不会因为多种子而翻号。厨房里同一天可以连续进门三块：一片摊开的面皮、一块肉饼、半只鸡胸。面板若给出同一条“组合、中等、25 分钟”，三块面对的是同一开环时钟。

\section{宽格子的分位}

宽格子在 648 个因子组合上各积大量种子，只保留分位，不把这张表并进 S1 的 19440。放大使用与 19440 行相同的差分和相同的种子相位 \(1000+\mathrm{seed}+17\times\mathrm{proto}\)，步长 \(2\,\mathrm{s}\)，记录间隔 \(10\,\mathrm{s}\)，时域 55 min。每个格子 30000 个种子，分位按排序后的线性插值，不是最近秩。图\ref{fig:wideband}给出对偶策略下停机偏差随厚度的分位带。15 mm 的带贴着上沿 \(+0.150\)；30 mm 的带跨过零点；50 mm 的带全部在零下，上沿仍低于 \(-0.3\)。带的宽度来自种子，厚度把整条带平移。把 30 个种子换成更宽的种子集，厚度的符号关系保持。图\ref{fig:s1wide}是宽格子的另一张总览：同一厚度上的竖向散开来自初温、含水和非均质，不是来自种子方差。

肉糜、30 mm、初温 \(5^\circ\mathrm{C}\)、含水 0.72 上，30000 个种子把 S2 的 40 个种子重复了出来：均匀对偶 \(e_{50}=-0.1113\)，提前比例 1；导热缺损对偶 \(-0.353\)，提前比例 1；均匀无探针 \(-0.5786\)；均匀开环 \(-0.8488\)；均匀腔温比例 \(+0.150\)，提前比例 0。种子从 40 增到 30000，没有把这些中位数移过零线。三种食品在 216 个设定里有 198 个中位数完全相同。其余 18 个的食品间极差最大为 0.040。原型进入的是噪声相位，进度律本身不区分食品。

肉糜 18 个初温、含水、非均质组合上，格中位数再取中位：15 mm 开环和腔温比例都是 \(+0.150\)，无探针 \(-0.366\)，对偶 \(-0.3387\)；30 mm 开环 \(-0.8382\)，腔温比例 \(+0.150\)，无探针 \(-0.4893\)，对偶 \(-0.0129\)；50 mm 开环 \(-0.8498\)，腔温比例 \(-0.2437\)，无探针 \(-0.4787\)，对偶 \(-0.4388\)。对偶格中位数大于零的只有两档，而且三种食品相同。一档是 15 mm、初温 \(5^\circ\mathrm{C}\)、导热缺损，三种含水的中位数都是 \(+0.150\)，十分位 \(-0.5255\)，提前比例 0.305：中位数落在进度上沿，约三成种子仍在设计线之下。另一档是 30 mm、初温 \(20^\circ\mathrm{C}\)、初始温度倾斜，中位数 \(+0.0945\)，十分位 \(+0.0733\)，提前比例 0。30 mm 对偶的 18 个格中位数本身从 \(-0.353\) 到 \(+0.0945\)，所以 \(-0.0129\) 不能代替按非均质分开的句子。30000 个种子收紧的是同一条确定性噪声相位上的分位，不是独立烹饪实验的置信区间。提前比例为 1 的格子，再增加种子也不会凭空产生跨过设计线的轨迹，除非种子相位本身改变停机。

\begin{figure}[htbp]
\centering
\includegraphics[width=0.86\textwidth]{S1_wide.png}
\caption{宽格子停机偏差总览。50 mm 的带不越过零。}
\label{fig:s1wide}
\end{figure}

\begin{figure}[htbp]
\centering
\includegraphics[width=0.86\textwidth]{X_wide_band.png}
\caption{宽格子、对偶策略：停机偏差随厚度的分位带。50 mm 的带不越过零。}
\label{fig:wideband}
\end{figure}

图\ref{fig:wideheat}是同一张表上的累计加热。腔温比例的带退化成一条水平线 3310，因为它不停机。开环的带按厚度菜单分成三档。对偶和无探针的带随厚度上升，厚件把加热器开到时域终点附近。图\ref{fig:wideunder}给出提前比例。无探针在 30 mm 与 50 mm 上接近 1；对偶在 50 mm 上同样是 1；开环在 15 mm 上提前比例为 0，因为菜单把薄件推过设计线。

\begin{figure}[htbp]
\centering
\includegraphics[width=0.82\textwidth]{X_wide_heat.png}
\caption{宽格子：累计加热的分位。腔温比例是一条水平线。}
\label{fig:wideheat}
\end{figure}

\begin{figure}[htbp]
\centering
\includegraphics[width=0.82\textwidth]{X_wide_under.png}
\caption{宽格子：停机进度低于设计线的比例。50 mm 上对偶与无探针都是 1。}
\label{fig:wideunder}
\end{figure}

S1 的小提琴图把四个策略的停机偏差摊开，见图\ref{fig:violin}。开环在 \(-0.85\) 形成高峰，薄件的上沿是另一峰。腔温比例的质量堆在 \(+0.150\)。无探针和对偶都在负半轴形成宽峰，对偶的峰更靠近零。经验分布见图\ref{fig:ecdf}：对偶的曲线在无探针右侧、开环左侧，仍然有九成样本落在零下。开环在 \(-0.85\) 形成台阶，腔温比例在 \(+0.150\) 形成台阶。密度图给出同一事实：开环高峰在 \(-0.84\)，腔温比例高峰在 \(+0.15\)。

\begin{figure}[htbp]
\centering
\includegraphics[width=0.84\textwidth]{X_s1_violin.png}
\caption{S1 四个策略停机偏差的小提琴图。开环双峰；腔温比例堆在上沿。}
\label{fig:violin}
\end{figure}

\begin{figure}[htbp]
\centering
\includegraphics[width=0.82\textwidth]{X_s1_ecdf.png}
\caption{S1 停机偏差的经验分布。对偶在无探针右侧，九成仍在零下。}
\label{fig:ecdf}
\end{figure}

能量与进度的帕累托图见图\ref{fig:pareto}。腔温比例占据右上角：能量最大，进度停在上沿。开环占据左下角：能量最小，进度最负。对偶落在中间偏左，无探针中间偏右。四个策略各自停在能量--进度平面的不同位置。厚件把对偶点拉向左下，薄件拉向右上。帕累托图上的点是格子中位。

\begin{figure}[htbp]
\centering
\includegraphics[width=0.82\textwidth]{X_s1_pareto.png}
\caption{S1：累计加热中位对停机偏差中位。对偶落在开环与腔温比例之间。}
\label{fig:pareto}
\end{figure}

图\ref{fig:underbar}按策略给出提前比例。无探针 0.944，对偶 0.902，开环 0.667，腔温比例 0.222。腔温比例的低提前比例来自不停机、进度饱和：它闭合的是空气温度。

\begin{figure}[htbp]
\centering
\includegraphics[width=0.78\textwidth]{X_s1_under.png}
\caption{S1 四个策略的提前比例 \(P(e_{\alpha}<0)\)。}
\label{fig:underbar}
\end{figure}

\section{集合散布}

集合积分把同一核的 32 或 128 个成员从同一初值推到时域终点，成员之间没有分析步、没有观测更新，只有散列噪声。表\ref{tab:ens}给出结束时的中位、散布和热阻方差。

\begin{table}[htbp]
\centering
\caption{同一核的成员散布。\(\mathrm{spread}_{\max}\) 是成员间停机偏差的极差。}
\label{tab:ens}
\begin{tabular}{rrrrrr}
\toprule
成员 & \(L/\mathrm{mm}\) & 非均质 & \(e_{50}\) & 极差 & 结束方差 \\
\midrule
32 & 15 & 0 & \(-0.109\) & 0.001 & 0.224 \\
32 & 15 & 2 & \(-0.461\) & 0.001 & 0.224 \\
32 & 30 & 0 & \(-0.096\) & 0.000 & 0.224 \\
32 & 30 & 2 & \(-0.203\) & 0.001 & 0.224 \\
32 & 50 & 0 & \(-0.707\) & 0.000 & 0.224 \\
32 & 50 & 2 & \(-0.802\) & 0.000 & 0.224 \\
128 & 50 & 0 & \(-0.707\) & 0.000 & 0.224 \\
128 & 50 & 2 & \(-0.802\) & 0.000 & 0.224 \\
\bottomrule
\end{tabular}
\end{table}

50 mm 上极差是 0。成员走的是同一条确定性轨迹，散列噪声在厚件的扩散时间尺度上被积分掉。15 mm 上极差约 0.001，仍比厚度效应小两个数量级。把成员数从 32 加到 128，中位与极差都不动。图\ref{fig:ensspread}与图\ref{fig:ensq}给出散布随时间和分位。这张图回答的是“同一核会不会自己散开”，不是“策略之间差多少”。结束方差 0.224 来自对偶探针的乘子，集合本身没有更新热阻。图\ref{fig:ensraw}给出原始成员曲线：50 mm 的线叠成一条，15 mm 才能看出一条细缝。若把集合散布误读成后验不确定度，会把已经看见的热阻重新写成噪声。

厨房里对应的判断是：同一台机器、同一块 50 mm 肉糜、同一套四个旋钮，把噪声相位换三十次，冷点仍然走不完导热时间尺度。面板上的“智能波动”若被理解成“多种子会给出不同的停机符号”，这张表给出的答案是不会。厚度把整条带平移，种子只在薄件上留下一条细缝。把 50 mm 时域从 55 min 放到 70 min，同一初值才走到设计线。时域不够，不是噪声不够。

\begin{figure}[htbp]
\centering
\includegraphics[width=0.82\textwidth]{E_ensemble.png}
\caption{同一核成员的进度曲线。50 mm 上成员重合，极差为 0。}
\label{fig:ensraw}
\end{figure}

\begin{figure}[htbp]
\centering
\includegraphics[width=0.82\textwidth]{X_ens_spread.png}
\caption{集合成员停机偏差的极差。50 mm 上极差为 0。}
\label{fig:ensspread}
\end{figure}

\begin{figure}[htbp]
\centering
\includegraphics[width=0.82\textwidth]{X_ens_quantile.png}
\caption{集合结束进度的分位。厚度平移整条分位，成员数从 32 到 128 不改变位置。}
\label{fig:ensq}
\end{figure}

\section{50 mm 的时域延长}

种子 7、肉糜 50 mm、组合档、对偶，55 min 时域上停机偏差 \(-0.782\)，停止时刻就是时域终点 3300 s，终点进度 0.068，累计加热 3310。把时域放到 70 min，同一初值在 4200 s 停机，偏差 \(+0.150\)，终点进度 0.9998，累计加热 3866。图\ref{fig:l50}给出两条时域的进度曲线。55 min 截在进度刚离开零的阶段；70 min 才走到设计线之上。

\begin{figure}[htbp]
\centering
\includegraphics[width=0.84\textwidth]{X_L50_horizon.png}
\caption{种子 7、50 mm、对偶：55 min 时域停在 \(e_{\alpha}=-0.782\)；70 min 时域走到上沿。}
\label{fig:l50}
\end{figure}

50 mm 的半厚 Fourier 数在 3300 s 为 \(\alpha t/L^2=2.0\times 10^{-7}\times 3300/0.05^2=0.264\)。导热时间尺度还没有走完。4200 s 时 \(\mathrm{Fo}=0.336\)，加上界面通量把表面维持在高温，内部才被推过 \(70^\circ\mathrm{C}\)。探针在 192 s 已经把方差降到 0.224，两条时域的结束方差相同。厚件到不了设计线，原因是扩散时间：热阻在 192 s 已经被看见。档位研究里 50 mm、对偶、组合的中位 \(e_{\alpha}=-0.782\)，与这条种子 7 轨迹同号同量级。延长时域是对中位的机理解释，主表仍按 55 min 报告。

厨房里对应的操作是：同一块厚肉糜，铃声在 55 min 响起之后，要不要再煮一刻钟。55 min 截在进度刚离开零的阶段，终点进度 0.068，\(e_{\alpha}=-0.782\)，累计加热 3310。再给 15 min，同一初值在 4200 s 走到上沿，进度 0.9998，累计加热 3866，结束方差仍是 0.224。多花的是扩散时间，不是再买一次热阻。水箱容量、门封和电气温升按 55 min 量级设计，家用过程很少被允许走到 70 min。面板若把 50 mm 仍写成“组合、中等、25 分钟”，给出的是开环菜单 1530 s，冷点更走不完。对偶把停机推到时域终点 3298 s，已经用满了家用上限，进度仍停在 0.068。集合积分把同一核的成员从 32 加到 128，50 mm 上极差仍为 0：噪声相位换三十次，冷点仍然走不完导热时间尺度。把“智能波动”理解成多种子会给出不同的停机符号，这张表给出的答案是不会。

水箱把这条上限写进了硬件。侧舱或底部水箱常见容量约 1 L，蒸汽阀拉满时每分钟送出数十克。55 min 组合档若一直维持基线蒸汽 0.15，耗水仍在设计范围内；探针窗把蒸汽拉到 1.0 只占用 40 s，额外消耗有限。把时域放到 70 min，胶条、管壁温升和水箱都开始离开标定区间。厚件缺的是扩散时间。说明书若把 50 mm 仍印成“中等、25 分钟”，给出的是开环 1530 s，冷点更走不完。对偶已经把加热用到 3298 s，进度仍是 0.068。宽格子里 50 mm 的上沿低于 \(-0.3\)，30000 个种子没有把这条带翻过零。厨房里多煮一刻钟，对应的是把积分终点从 55 min 挪到 70 min，不是再喷一轮蒸汽。

\section{档位研究上的初温与含水}

图\ref{fig:t0m0}把档位研究里对偶策略的停机偏差按初温和初始含水分面。初温从 \(4^\circ\mathrm{C}\) 到 \(20^\circ\mathrm{C}\) 把 30 mm 的中位上移；含水三档在同一厚度上几乎重合，与 S1 的零效应一致。图\ref{fig:studystop}给出停止时刻对累计加热。开环落在菜单时间的三条竖线上；腔温比例落在 3300 s 的竖线上；对偶和无探针填满中间。

\begin{figure}[htbp]
\centering
\includegraphics[width=0.88\textwidth]{X_study_t0m0.png}
\caption{档位研究、对偶：初温与初始含水的分面。含水三档重合。}
\label{fig:t0m0}
\end{figure}

\begin{figure}[htbp]
\centering
\includegraphics[width=0.82\textwidth]{X_study_stop_energy.png}
\caption{档位研究：停止时刻对累计加热。开环是三条菜单竖线。}
\label{fig:studystop}
\end{figure}

图\ref{fig:facet}把策略与档位、厚度画成分面。蒸档一行全部偏负；烤档与组合档只在 15 mm 与部分 30 mm 上为正；50 mm 一列全部为负。宽格子里厚度把整条分位带平移，50 mm 的上沿仍低于 \(-0.3\)；同一核的成员在 50 mm 上极差为 0；把时域放到 70 min，同一初值才走到上沿，结束方差仍是 0.224。策略名称写在允许集合里，几何先于菜单。

家用说明书把 55 min 写成“最长工作时间”或根本不写，只给出各菜单的分钟数。水箱约 1 L，蒸汽阀拉满时每分钟送出数十克。组合档基线蒸汽 0.15 开满 55 min，耗水仍在设计范围内；探针窗把蒸汽拉到 1.0 只占用 40 s，额外消耗有限。把同一初值放到 70 min，胶条、管壁温升和水箱都开始离开标定区间。种子 7、50 mm、对偶、组合在 55 min 停在进度 0.068，\(e_{\alpha}=-0.782\)，累计加热 3310；70 min 走到上沿，进度 0.9998，累计加热 3866，结束方差仍是 0.224。多花的 15 min 是扩散时间。厨房里若在 55 min 铃声后再煮一刻钟，对应的是把积分终点挪到 70 min，不是再喷一轮蒸汽。宽格子里 50 mm 的上沿低于 \(-0.3\)，30000 个种子没有把这条带翻过零。集合积分把成员从 32 加到 128，50 mm 上极差仍为 0。噪声相位换三十次，冷点仍然走不完导热时间尺度。

面板上的“智能波动”若被理解成多种子会给出不同的停机符号，这张表给出的答案是不会。厚度把整条带平移，种子只在薄件上留下一条细缝。15 mm 上极差约 0.001，仍比厚度效应小两个数量级。操作者同一天里进门三次，三次都是 50 mm 厚肉、同一套四个旋钮，铃声可以响在时域终点，切开之后中心仍停在粉色附近。把第三次换成 15 mm 面皮，开环菜单 900 s 已经把进度送到上沿。机器没有换，几何换了。分面图把这件事画成一列：50 mm 三种档位全部为负，蒸档一行全部偏负，烤档与组合档只在薄件和部分 30 mm 上为正。读这张图时先沿厚度走，再比较档位名称。

\begin{figure}[htbp]
\centering
\includegraphics[width=0.90\textwidth]{X_study_facet.png}
\caption{档位研究：策略、档位与厚度的分面。50 mm 一列全部低于设计线。}
\label{fig:facet}
\end{figure}
